% docs/translation/latex/front/translation_reading.tex - front matter of the translation report (PDF only): how to read this book.
% Hand-written LaTeX (no Markdown counterpart); \input by the main .tex of the book.
%   Compile: cmake --build <build> --target pdf, or python3 -I tools/build_docs_pdf.py --source-root . --build-dir DIR
%   (figures, then xelatex passes + makeindex -s docs/common/latex/simcpp.ist); by hand see docs/README.md.
%   XeLaTeX; class docs/common/latex/simcpp.cls (TeX Live packages listed there); fonts TeX Gyre Pagella,
%   TeX Gyre Heros, DejaVu Sans Mono. The figures are built first from docs/{manual,translation}/figures/*.tex
%   (preview.sty) into PDFs.
\section*{How to read this report}
The page \hyperref[dashboard]{\emph{The translation at a glance}} summarises the numbers. \textbf{Method} explains
how the Fortran program was translated and how each piece was proved against it, from single routines to complete
runs. \textbf{The mapping} is the reference for maintainers: which C++ code replaces which Fortran routine, the
numerical rules that give bit-for-bit identical results, and the switch \code{-{}-fortran-compat}. \textbf{Findings}
is the catalogue of every Fortran bug, quirk and deviation, with what SIMC++ does in each mode. The appendices hold the
generated table of all Fortran routines and the two user guides about the differences.
